%=============================================================================
% Chapter 2 Exercises
%=============================================================================
\newpage
\section{Chapter 2 Exercises [430]}

%-----------------------------------------------------------------------------
\subsection{[5] C to MIPS: Arithmetic}
\hypertarget{prob-2.1}{}
\noindent$\langle$\S2.2$\rangle$\medskip

\noindent For the following C statement, what is the corresponding MIPS assembly
code? Assume that the variables \texttt{f}, \texttt{g}, \texttt{h} have
already been placed in registers \texttt{\$s0}, \texttt{\$s1}, and
\texttt{\$s2}, respectively. Use a minimal number of MIPS assembly
instructions.
\begin{verbatim}
    f = g + h - 5;
\end{verbatim}

\problembreak

%-----------------------------------------------------------------------------
\subsection{[5] MIPS to C: Two Adds}
\hypertarget{prob-2.2}{}
\noindent$\langle$\S2.2$\rangle$\medskip

\noindent Write a single C statement that corresponds to the two MIPS assembly
instructions below.
\begin{verbatim}
    add $t0, $s1, $s2
    add $s0, $t0, $s3
\end{verbatim}

\problembreak

%-----------------------------------------------------------------------------
\subsection{[5] C to MIPS: Array Access}
\hypertarget{prob-2.3}{}
\noindent$\langle$\S\S2.2, 2.3$\rangle$\medskip

\noindent For the following C statement, write the corresponding MIPS assembly
code. Assume that the variables \texttt{f}, \texttt{g}, \texttt{h},
\texttt{i}, and \texttt{j} are assigned to registers \texttt{\$s0},
\texttt{\$s1}, \texttt{\$s2}, \texttt{\$s3}, and \texttt{\$s4},
respectively. Assume that the base address of the arrays \texttt{A} and
\texttt{B} are in registers \texttt{\$s6} and \texttt{\$s7}, respectively.
\begin{verbatim}
    B[8] = A[i-j];
\end{verbatim}

\problembreak

%-----------------------------------------------------------------------------
\subsection{[5] MIPS to C: Array Code}
\hypertarget{prob-2.4}{}
\noindent$\langle$\S2.2, 2.5$\rangle$\medskip

\noindent Translate the following MIPS code to C. Assume that the variables
\texttt{f}, \texttt{g}, \texttt{h}, \texttt{i}, and \texttt{j} are assigned
to registers \texttt{\$s0} through \texttt{\$s4}. Assume that the base
addresses of the arrays \texttt{A} and \texttt{B} are in registers
\texttt{\$s6} and \texttt{\$s7}, respectively.
\begin{verbatim}
    sll  $t0, $s0, 2       # $t0 = f * 4
    add  $t0, $s6, $t0     # $t0 = &A[f]
    sll  $t1, $s1, 2       # $t1 = g * 4
    add  $t1, $s7, $t1     # $t1 = &B[g]
    lw   $s0, 0($t0)       # f = A[f]
    addi $t2, $t0, 4
    lw   $t0, 0($t2)
    add  $t0, $t0, $s0
    sw   $t0, 0($t1)
\end{verbatim}

\problembreak

%-----------------------------------------------------------------------------
\subsection{[5] Byte Ordering: Endianness}
\hypertarget{prob-2.5}{}
\noindent$\langle$\S2.3$\rangle$\medskip

\noindent Show how the value \texttt{0xabcdef12} would be arranged in memory
of a little-endian and a big-endian machine. Assume that the data are stored
starting at address 0 and that the word size is 4 bytes.

\problembreak

%-----------------------------------------------------------------------------
\subsection{[5] Hex to Decimal Conversion}
\hypertarget{prob-2.6}{}
\noindent$\langle$\S2.4$\rangle$\medskip

\noindent Translate \texttt{0xabcdef12} into decimal.

\problembreak

%-----------------------------------------------------------------------------
\subsection{[5] C to MIPS: 8-Byte Word Arrays}
\hypertarget{prob-2.7}{}
\noindent$\langle$\S2.2, 2.3$\rangle$\medskip

\noindent Translate the following C code to MIPS. Assume that the variables
\texttt{f}, \texttt{g}, \texttt{h}, \texttt{i}, and \texttt{j} are assigned
to registers \texttt{\$s0} through \texttt{\$s4}. Assume that the base
addresses of the arrays \texttt{A} and \texttt{B} are in registers
\texttt{\$s6} and \texttt{\$s7}, respectively. Assume that the elements of
the arrays \texttt{A} and \texttt{B} are 8-byte words.
\begin{verbatim}
    B[8] = A[i] + A[j];
\end{verbatim}

\problembreak

%-----------------------------------------------------------------------------
\subsection{[10] Instruction Format Fields}
\hypertarget{prob-2.8}{}
\noindent$\langle$\S2.5$\rangle$\medskip

\noindent For each MIPS instruction below, show the value of the opcode
(\texttt{op}), source register (\texttt{rs}), funct field, and destination
register (\texttt{rd}) fields. For I-type instructions, show the value of the
immediate field. For R-type instructions, show the value of the second source
register (\texttt{rt}).
\begin{verbatim}
    add $s0, $s0, $s1
    sw  $t1, 0($s0)
\end{verbatim}

\problembreak

%-----------------------------------------------------------------------------
\subsection{[20] MIPS Fields for Exercise 2.8}
\hypertarget{prob-2.9}{}
\noindent$\langle$\S2.3, 2.5$\rangle$\medskip

\noindent For each MIPS instruction in Exercise 2.8, show the value of the
opcode (op), source register (rs) and funct field, and destination register
(rd) fields. For I-type instructions, show the value of the immediate field,
and for R-type instructions, show the value of the second source register (rt).

\problembreak

%-----------------------------------------------------------------------------
\subsection{[30] Register Arithmetic and Overflow}
\hypertarget{prob-2.10}{}
\noindent$\langle$\S2.4$\rangle$\medskip

\noindent Assume that registers \texttt{\$s0} and \texttt{\$s1} hold the
values \texttt{0x80000000} and \texttt{0xD0000000}, respectively.

\begin{enumerate}
  \phantomsection\addcontentsline{toc}{subsubsection}{2.10.1~[5]}
  \item[2.10.1] [5] What is the value of \texttt{\$t0} for the following assembly
        code? \quad \texttt{add \$t0, \$s0, \$s1}
  \phantomsection\addcontentsline{toc}{subsubsection}{2.10.2~[5]}
  \item[2.10.2] [5] Is the result in \texttt{\$t0} the desired result, or has
        there been overflow?
  \phantomsection\addcontentsline{toc}{subsubsection}{2.10.3~[5]}
  \item[2.10.3] [5] For the contents of registers \texttt{\$s0} and
        \texttt{\$s1} as specified above, what is the value of \texttt{\$t0}
        for the following assembly code? \quad \texttt{sub \$t0, \$s0, \$s1}
  \phantomsection\addcontentsline{toc}{subsubsection}{2.10.4~[5]}
  \item[2.10.4] [5] Is the result in \texttt{\$t0} the desired result, or has
        there been overflow?
  \phantomsection\addcontentsline{toc}{subsubsection}{2.10.5~[5]}
  \item[2.10.5] [5] For the contents of registers \texttt{\$s0} and
        \texttt{\$s1} as specified above, what is the value of \texttt{\$t0}
        for the following assembly code?
        \quad \texttt{add \$t0, \$s0, \$s1} \quad followed by
        \quad \texttt{add \$t0, \$t0, \$s0}
  \phantomsection\addcontentsline{toc}{subsubsection}{2.10.6~[5]}
  \item[2.10.6] [5] Is the result in \texttt{\$t0} the desired result, or has
        there been overflow?
\end{enumerate}

\problembreak

%-----------------------------------------------------------------------------
\subsection{[15] Overflow Ranges}
\hypertarget{prob-2.11}{}
\noindent$\langle$\S2.4$\rangle$\medskip

\noindent Assume that \texttt{\$s1} holds the value $128_{10}$.

\begin{enumerate}
  \phantomsection\addcontentsline{toc}{subsubsection}{2.11.1~[5]}
  \item[2.11.1] [5] For the instruction \texttt{add \$s0, \$s0, \$s1}, what is
        the range(s) of values for \texttt{\$s0} that would result in
        overflow?
  \phantomsection\addcontentsline{toc}{subsubsection}{2.11.2~[5]}
  \item[2.11.2] [5] For the instruction \texttt{sub \$s0, \$s0, \$s1}, what is
        the range(s) of values for \texttt{\$s0} that would result in
        overflow?
  \phantomsection\addcontentsline{toc}{subsubsection}{2.11.3~[5]}
  \item[2.11.3] [5] For the instruction \texttt{sub \$s0, \$s1, \$s0}, what is
        the range(s) of values for \texttt{\$s0} that would result in
        overflow?
\end{enumerate}

\problembreak

%-----------------------------------------------------------------------------
\subsection{[5] Binary to Instruction Type}
\hypertarget{prob-2.12}{}
\noindent$\langle$\S2.4, 2.5$\rangle$\medskip

\noindent Provide the type and assembly language instruction for the following
binary value:
\begin{verbatim}
    0000 0010 0001 0100 1000 0000 0010 0000
\end{verbatim}
Hint: Figure 2.20 may be helpful.

\problembreak

%-----------------------------------------------------------------------------
\subsection{[5] Instruction to Hex Representation}
\hypertarget{prob-2.13}{}
\noindent$\langle$\S2.4, 2.5$\rangle$\medskip

\noindent Provide the type and hexadecimal representation of the following
instructions:
\begin{verbatim}
    add $t1, $s2, $t1
\end{verbatim}

\problembreak

%-----------------------------------------------------------------------------
\subsection{[5] MIPS Fields to Instruction (R-type)}
\hypertarget{prob-2.14}{}
\noindent$\langle$\S2.5$\rangle$\medskip

\noindent Provide the type, assembly language instruction, and binary
representation of the instruction described by the following MIPS fields:
\texttt{op=0, rs=3, rt=2, rd=3, shamt=0, funct=34}.

\problembreak

%-----------------------------------------------------------------------------
\subsection{[5] MIPS Fields to Instruction (I-type)}
\hypertarget{prob-2.15}{}
\noindent$\langle$\S2.5$\rangle$\medskip

\noindent Provide the type, assembly language instruction, and binary
representation of the instruction described by the following MIPS fields:
\texttt{op=0x23, rs=1, rt=2, const=0x4}.

\problembreak

%-----------------------------------------------------------------------------
\subsection{[15] Expanding the Register File}
\hypertarget{prob-2.16}{}
\noindent$\langle$\S2.5$\rangle$\medskip

\noindent Assume that we would like to expand the MIPS register file to 128
registers and expand the instruction set to contain four times as many
instructions.

\begin{enumerate}
  \phantomsection\addcontentsline{toc}{subsubsection}{2.16.1~[5]}
  \item[2.16.1] [5] How would this affect the size of each of the bit fields in
        the R-type instructions?
  \phantomsection\addcontentsline{toc}{subsubsection}{2.16.2~[5]}
  \item[2.16.2] [5] How would this affect the size of each of the bit fields in
        the I-type instructions?
  \phantomsection\addcontentsline{toc}{subsubsection}{2.16.3~[5]}
  \item[2.16.3] [5] How could each of two proposed changes decrease the size of
        an MIPS assembly program? On the other hand, how could the proposed
        change increase the size of an MIPS assembly program?
\end{enumerate}

\problembreak

%-----------------------------------------------------------------------------
\subsection{[15] Bitwise Operations on Registers}
\hypertarget{prob-2.17}{}
\noindent$\langle$\S2.6$\rangle$\medskip

\noindent Assume the following register contents:
\texttt{\$t0 = 0xAAAAAAAA}, \texttt{\$t1 = 0x12345678}.

\begin{enumerate}
  \phantomsection\addcontentsline{toc}{subsubsection}{2.17.1~[5]}
  \item[2.17.1] [5] What is the value of \texttt{\$t2} after the following
        instructions?
\begin{verbatim}
    sll $t2, $t0, 4
    or  $t2, $t2, $t1
\end{verbatim}
  \phantomsection\addcontentsline{toc}{subsubsection}{2.17.2~[5]}
  \item[2.17.2] [5] What is the value of \texttt{\$t2} after the following
        instructions?
\begin{verbatim}
    srl $t2, $t0, 3
    and $t2, $t2, $t1
\end{verbatim}
  \phantomsection\addcontentsline{toc}{subsubsection}{2.17.3~[5]}
  \item[2.17.3] [5] What is the value of \texttt{\$t2} after the following
        instructions?
\begin{verbatim}
    srl $t2, $t0, 4
    andi $t2, $t2, -1
\end{verbatim}
\end{enumerate}

\problembreak

%-----------------------------------------------------------------------------
\subsection{[10] Bit Field Extraction}
\hypertarget{prob-2.18}{}
\noindent$\langle$\S2.6$\rangle$\medskip

\noindent Find the shortest sequence of MIPS instructions that extracts bits
16 down to 11 from register \texttt{\$t0} and uses the value of this field to
replace bits 31 down to 26 in register \texttt{\$t1} without changing the
other bits of registers \texttt{\$t0} and \texttt{\$t1}. (Be sure to test
your code using \texttt{\$t0 = 0} and \texttt{\$t1 = 0xFFFFFFFF}. Doing so
may reveal a common oversight.)

\problembreak

%-----------------------------------------------------------------------------
\subsection{[5] Pseudoinstruction Implementation}
\hypertarget{prob-2.19}{}
\noindent$\langle$\S2.6$\rangle$\medskip

\noindent Provide a minimal set of MIPS instructions that may be used to
implement the following pseudoinstruction:
\begin{verbatim}
    not $t0, $t1
\end{verbatim}

\problembreak

%-----------------------------------------------------------------------------
\subsection{[5] C to MIPS: Array Sum}
\hypertarget{prob-2.20}{}
\noindent$\langle$\S2.6$\rangle$\medskip

\noindent For the following C statement, write a minimal sequence of MIPS
assembly instructions that does the identical operation. Assume
\texttt{\$t0 = A} and \texttt{\$s0} is the base address of \texttt{C}.
\begin{verbatim}
    A = C[0] + C[4];
\end{verbatim}

\problembreak

%-----------------------------------------------------------------------------
\subsection{[5] SLT and Branch}
\hypertarget{prob-2.21}{}
\noindent$\langle$\S2.6$\rangle$\medskip

\noindent Assume \texttt{\$s0} holds the value \texttt{0x80000000}. What is
the value of \texttt{\$t2} after the following instructions?
\begin{verbatim}
    slt  $t2, $s0, $zero
    bne  $t2, $zero, ELSE
    j    EXIT
ELSE:
    addi $t2, $t2, 2
EXIT:
\end{verbatim}

\problembreak

%-----------------------------------------------------------------------------
\subsection{[10] PC Addressing Range}
\hypertarget{prob-2.22}{}
\noindent$\langle$\S2.10$\rangle$\medskip

\noindent Suppose the program counter (PC) is set to \texttt{0x00400000}.

\begin{enumerate}
  \phantomsection\addcontentsline{toc}{subsubsection}{2.22.1~[5]}
  \item[2.22.1] [5] What range of addresses can be reached using the MIPS
        jump-and-link (\texttt{jal}) instruction? (In other words, what is
        the set of possible values for the PC after the jump instruction
        executes?)
  \phantomsection\addcontentsline{toc}{subsubsection}{2.22.2~[5]}
  \item[2.22.2] [5] What range of addresses can be reached using the MIPS branch
        if equal (\texttt{beq}) instruction? (In other words, what is the set
        of possible values for the PC after the branch instruction executes?)
\end{enumerate}

\problembreak

%-----------------------------------------------------------------------------
\subsection{[5] Proposed \texttt{rpt} Instruction}
\hypertarget{prob-2.23}{}
\medskip

\noindent Consider a proposed new instruction named \texttt{rpt}. This
instruction combines a loop's condition check and counter decrement into a
single instruction. For example, \texttt{rpt \$s0, 1, LOOP} would do the
following:

\begin{verbatim}
    if ($s0 > 0)
        $s0 = $s0 - 1
        goto LOOP
\end{verbatim}

\noindent What instruction format would be a good fit for this instruction?
What are the issues that might prevent this from being a real MIPS instruction?

\problembreak

%-----------------------------------------------------------------------------
\subsection{[15] MIPS Loop Analysis}
\hypertarget{prob-2.24}{}
\noindent$\langle$\S2.7$\rangle$\medskip

\noindent Consider the following MIPS loop:
\begin{verbatim}
          add  $s2, $zero, $zero
    LOOP: lw   $s1, 0($s0)
          add  $s2, $s2, $s1
          addi $s0, $s0, 4
          addi $t1, $t1, -1
          bne  $t1, $zero, LOOP
\end{verbatim}

\begin{enumerate}
  \phantomsection\addcontentsline{toc}{subsubsection}{2.24.1~[5]}
  \item[2.24.1] [5] Assume that the register \texttt{\$t1} is initialized to the
        value 10. What is the value in register \texttt{\$s2} assuming
        \texttt{\$s0} initially points to an array of integers?
  \phantomsection\addcontentsline{toc}{subsubsection}{2.24.2~[5]}
  \item[2.24.2] [5] For the loop above, write the equivalent C code. Assume that
        the registers \texttt{\$s1}, \texttt{\$s2}, \texttt{\$t1}, and
        \texttt{\$t2} are integers \texttt{a}, \texttt{b}, \texttt{i}, and
        \texttt{temp}, respectively.
  \phantomsection\addcontentsline{toc}{subsubsection}{2.24.3~[5]}
  \item[2.24.3] [5] For the loop above, assume that \texttt{\$t1} is initialized
        to the value $N$. How many MIPS instructions are executed?
\end{enumerate}

\problembreak

%-----------------------------------------------------------------------------
\subsection{[10] C to MIPS: Nested Loop}
\hypertarget{prob-2.25}{}
\noindent$\langle$\S2.7$\rangle$\medskip

\noindent Translate the following C code to MIPS assembly code. Use a minimum
number of instructions. Assume that the values of \texttt{a}, \texttt{b},
\texttt{i}, and \texttt{j} are in registers \texttt{\$s0}, \texttt{\$s1},
\texttt{\$t0}, and \texttt{\$t1}, respectively. Also, assume that register
\texttt{\$s2} holds the base address of array \texttt{D}.
\begin{verbatim}
    for (i = 0; i < a; i++)
        for (j = 0; j < b; j++)
            D[4*j] = i + j;
\end{verbatim}

\problembreak

%-----------------------------------------------------------------------------
\subsection{[5] Instruction Count for Nested Loop}
\hypertarget{prob-2.26}{}
\noindent$\langle$\S2.7$\rangle$\medskip

\noindent How many MIPS instructions does it take to implement the C code from
Exercise 2.25? If the variables \texttt{a} and \texttt{b} are initialized to
10 and 1 and all elements of \texttt{D} are initially 0, what is the total
number of MIPS instructions that is executed to complete the loop?

\problembreak

%-----------------------------------------------------------------------------
\subsection{[5] MIPS Loop to C}
\hypertarget{prob-2.27}{}
\noindent$\langle$\S2.7$\rangle$\medskip

\noindent Translate the following loop into C. Assume that the C-level integer
\texttt{i} is held in register \texttt{\$s1}; \texttt{\$s2} holds the C-level
integer called \texttt{result}; and \texttt{\$s0} holds the base address of
the integer array \texttt{MemArray}.
\begin{verbatim}
          addi $s1, $zero, 0
    LOOP: lw   $s2, 0($s0)
          add  $s2, $s2, $s1
          addi $s0, $s0, 4
          addi $s1, $s1, 1
          slti $t0, $s1, 10
          bne  $t0, $zero, LOOP
\end{verbatim}

\problembreak

%-----------------------------------------------------------------------------
\subsection{[10] Optimize MIPS Loop}
\hypertarget{prob-2.28}{}
\noindent$\langle$\S2.7$\rangle$\medskip

\noindent Rewrite the loop from Exercise 2.27 to reduce the number of MIPS
instructions executed. Hint: Notice that variable \texttt{i} is used only for
loop control.

\problembreak

%-----------------------------------------------------------------------------
\subsection{[30] Fibonacci in MIPS}
\hypertarget{prob-2.29}{}
\noindent$\langle$\S2.8$\rangle$\medskip

\noindent Implement the following C code in MIPS assembly. Hint: Remember that
the stack pointer must remain aligned on a multiple of 16.
\begin{verbatim}
    int fib(int n) {
        if (n == 0) return 0;
        else if (n == 1) return 1;
        else return fib(n-1) + fib(n-2);
    }
\end{verbatim}

\problembreak

%-----------------------------------------------------------------------------
\subsection{[20] Stack Contents for Function Calls}
\hypertarget{prob-2.30}{}
\noindent$\langle$\S2.8$\rangle$\medskip

\noindent For each function call, show the contents of the stack after the
function call is made. Assume the stack pointer is initially at address
\texttt{0x7FFFFFFC}, and follow the register conventions as specified in
Figure 2.11.

\problembreak

%-----------------------------------------------------------------------------
\subsection{[20] Translate Function to MIPS}
\hypertarget{prob-2.31}{}
\noindent$\langle$\S2.8$\rangle$\medskip

\noindent Translate function \texttt{f} into MIPS assembly language. If you
need to use registers \texttt{\$t5} through \texttt{\$t7}, use the
lower-numbered registers first. Assume the function declaration for
\texttt{func} is \texttt{int func(int a, int b);}. The code for function
\texttt{f} is as follows:
\begin{verbatim}
    int f(int a, int b, int c, int d) {
        return func(func(a, b), c + d);
    }
\end{verbatim}

\problembreak

%-----------------------------------------------------------------------------
\subsection{[5] Tail-Call Optimization}
\hypertarget{prob-2.32}{}
\noindent$\langle$\S2.8$\rangle$\medskip

\noindent Can we use the tail-call optimization in the function from Exercise
2.31? If not, explain why not. If yes, what is the difference in the number of
instructions executed with and without the optimization?

\problembreak

%-----------------------------------------------------------------------------
\subsection{[5] Stack State Before Return}
\hypertarget{prob-2.33}{}
\noindent$\langle$\S2.8$\rangle$\medskip

\noindent Right before your function \texttt{f} from Exercise 2.31 returns,
what do we know about contents of registers \texttt{\$s1}, \texttt{\$s3},
\texttt{\$s4}, and \texttt{\$t5}? Keep in mind that we know what the entire
function \texttt{f} looks like, but for \texttt{func} we only know its
declaration.

\problembreak

%-----------------------------------------------------------------------------
\subsection{[30] ASCII String to Integer}
\hypertarget{prob-2.34}{}
\noindent$\langle$\S2.9$\rangle$\medskip

\noindent Write a program in MIPS assembly language to convert an ASCII number
string containing positive and negative integer decimal strings, to an
integer. Your program should expect register \texttt{\$a0} to hold the
address of a null-terminated string containing some combination of the digits
0 through 9. Your program should compute the integer value equivalent to this
string of digits, then place the number in register \texttt{\$v0}. If a
non-digit character appears anywhere in the string, your program should stop
with the value $-1$ in register \texttt{\$v0}. For example, if register
\texttt{\$a0} points to a sequence of three bytes
$50_{10}$, $52_{10}$, $0_{10}$ (the null-terminated string ``24''), then when
the program stops, register \texttt{\$v0} should contain the value $24_{10}$.
The RISC-V null instruction takes two registers as input.
There is no ``null'' instruction. Just store the constant 10 in a register.

\problembreak

%-----------------------------------------------------------------------------
\subsection{[10] Memory Store and Endianness}
\hypertarget{prob-2.35}{}
\noindent$\langle$\S2.9$\rangle$\medskip

\noindent For the following code:
\begin{verbatim}
    lbu $t0, 0($t1)
    sw  $t0, 0($t1)
\end{verbatim}

\noindent Assume that the register \texttt{\$t1} contains the address
\texttt{0x10000000}.

\begin{enumerate}
  \phantomsection\addcontentsline{toc}{subsubsection}{2.35.1~[5]}
  \item[2.35.1] [5] What value is stored at \texttt{0x10000000} on a big-endian
        machine?
  \phantomsection\addcontentsline{toc}{subsubsection}{2.35.2~[5]}
  \item[2.35.2] [5] What value is stored at \texttt{0x10000000} on a
        little-endian machine?
\end{enumerate}

\problembreak

%-----------------------------------------------------------------------------
\subsection{[5] Create 32-bit Constant}
\hypertarget{prob-2.36}{}
\noindent$\langle$\S2.10$\rangle$\medskip

\noindent Write the MIPS assembly code that creates the 32-bit constant
\texttt{0x0000FACE} and stores that value to register \texttt{\$t1}.

\problembreak

%-----------------------------------------------------------------------------
\subsection{[10] Atomic Set-Max with \texttt{ll/sc}}
\hypertarget{prob-2.37}{}
\noindent$\langle$\S2.11$\rangle$\medskip

\noindent Write the MIPS assembly code to implement the following C code as an
atomic ``set max'' operation using the \texttt{ll} and \texttt{sc}
instructions. Hint: the argument \texttt{x} contains the address of a shared
variable, which should be replaced by \texttt{s} if \texttt{s} is greater
than the value it points to:
\begin{verbatim}
    void setmax(int* x, int s) {
        // Begin critical section
        if (s > *x)
            *x = s;
        // End of critical section
    }
\end{verbatim}

\problembreak

%-----------------------------------------------------------------------------
\subsection{[5] Concurrent Critical Section}
\hypertarget{prob-2.38}{}
\noindent$\langle$\S2.11$\rangle$\medskip

\noindent Using your code from Exercise 2.37 as an example, explain what
happens when two processors begin to execute the critical section at the same
time, assuming that each processor executes exactly one instruction per cycle.

\problembreak

%-----------------------------------------------------------------------------
\subsection{[10] New Arithmetic Instructions}
\hypertarget{prob-2.39}{}
\noindent$\langle$\S2.14, 2.13$\rangle$\medskip

\noindent Assume for a given processor the CPI of arithmetic instructions is
1, the CPI of load/store instructions is 10, and the CPI of branch
instructions is 3. Assume a program has the following instruction breakdown:
500 million arithmetic instructions, 300 million load/store instructions,
100 million branch instructions.

\begin{enumerate}
  \phantomsection\addcontentsline{toc}{subsubsection}{2.39.1~[5]}
  \item[2.39.1] [5] Suppose that new, more powerful arithmetic instructions are
        added to the instruction set. On average, through the use of these
        more powerful arithmetic instructions, we can reduce the number of
        arithmetic instructions needed to execute a program by 25\%, and the
        cost of increasing the clock cycle time by 10\%. Is this a good
        design choice? Why?
  \phantomsection\addcontentsline{toc}{subsubsection}{2.39.2~[5]}
  \item[2.39.2] [5] Suppose that we find a way to double the performance of
        arithmetic instructions. What is the overall speedup of our machine?
        What if we find a way to improve the performance of arithmetic
        instructions by 10 times?
\end{enumerate}

\problembreak

%-----------------------------------------------------------------------------
\subsection{[15] CPI and Instruction Mix}
\hypertarget{prob-2.40}{}
\noindent$\langle$\S2.14, 2.13$\rangle$\medskip

\noindent Assume that for a given program 70\% of the executed instructions
are arithmetic, 10\% are load/store, and 20\% are branch.

\begin{enumerate}
  \phantomsection\addcontentsline{toc}{subsubsection}{2.40.1~[5]}
  \item[2.40.1] [5] Given the instruction mix and the assumption that an
        arithmetic instruction requires 2 cycles, a load/store instruction
        takes 6 cycles, and a branch instruction takes 3 cycles, find the
        average CPI.
  \phantomsection\addcontentsline{toc}{subsubsection}{2.40.2~[5]}
  \item[2.40.2] [5] For a 25\% improvement in performance, how many cycles, on
        average, may an arithmetic instruction take if load/store and branch
        instructions are not improved at all?
  \phantomsection\addcontentsline{toc}{subsubsection}{2.40.3~[5]}
  \item[2.40.3] [5] For a 50\% improvement in performance, how many cycles, on
        average, may an arithmetic instruction take if load/store and branch
        instructions are not improved at all?
\end{enumerate}

\problembreak

%-----------------------------------------------------------------------------
\subsection{[10] Scaled Offset for Exercise 2.4}
\hypertarget{prob-2.41}{}
\noindent$\langle$\S2.21$\rangle$\medskip

\noindent Suppose the MIPS ISA included a scaled offset addressing mode
similar to the x86 one described in Section 2.19 (Figure 2.35). Describe how
you would use scaled offset loads to further reduce the number of assembly
instructions needed to carry out the function given in Exercise 2.4.

\problembreak

%-----------------------------------------------------------------------------
\subsection{[10] Scaled Offset for Exercise 2.7}
\hypertarget{prob-2.42}{}
\noindent$\langle$\S2.21$\rangle$\medskip

\noindent Suppose the MIPS ISA included a scaled offset addressing mode
similar to the x86 one described in Section 2.19 (Figure 2.35). Describe how
you would use scaled offset loads to further reduce the number of assembly
instructions needed to implement the C code given in Exercise 2.7.
